% == == == == == == == == == == == == == == == == == == == == == == == == == ==
    == == == == == == == == == == == ==
    = % NASA POWER OF 10 RULES FOR SAFETY - CRITICAL CODE % == == == == == == ==
      == == == == == == == == == == == == == == == == == == == == == == == == ==
      == == == == == ==
    =

\subsection{NASA Power of 10: Rules for Safety-Critical Code
}
\label{sec:nasa-power-of-10}

This section documents the NASA/JPL Power of 10 coding rules developed by Gerard J. Holzmann in 2006 for safety-critical embedded systems. These rules prioritize verification and predictability over coding flexibility.

\textbf{STAR Compliance Philosophy:} The STAR project follows the Power of 10 rules, with one \textbf{intentional architectural deviation
} for Dependency Inversion Principle (DIP) -- a conscious trade-off for testability and modularity.

\subsubsection{Compliance Status Legend
 }

\begin{table}[H]
\centering
\begin{tabular} { | l | l |}
\hline
\textbf{Symbol} & \textbf{Meaning}
\\
\hline
\textcolor{green} {\textbf{$\checkmark$ COMPLIANT}} & STAR follows this rule \\
\hline
\textcolor{orange} {\textbf{$\triangle$ !INTENTIONAL DEVIATION}} &
    Architectural choice(DIP pattern) \\
\hline
\end{tabular}
\end{table}

\subsubsection{Rule 1 : Simplify Control Flow}

\textbf{Rule:} Restrict all code to very simple control flow constructs -- no \texttt{goto}, \texttt{setjmp}/\texttt{longjmp}, or recursion (direct or indirect).

\textbf{Rationale:} Creates predictable program structure for human understanding and static analysis. Recursion causes unbounded stack usage -- dangerous in embedded systems.

\textbf{STAR Compliance:
}
\textcolor{green} {
  \textbf{$\checkmark$ COMPLIANT}
}

\begin{lstlisting}[language = C]
    /* COMPLIANT - No goto, no recursion */
    esp_err_t
    star_motor_init(star_motor_handle_t *handle,
                    const star_motor_config_t *config) {
  if (handle == NULL || config == NULL) {
    return ESP_ERR_INVALID_ARG;
 }

  /* All control flow uses if/while/for */
  esp_err_t ret = configure_timer(handle, config);
  if (ret != ESP_OK) {
    return ret;
 }

  return configure_gpio(handle, config);
}
\end{lstlisting}

\subsubsection{Rule 2 : Fixed Loop Upper - Bounds}

\textbf{Rule :} All loops must have a fixed upper bound
    that static analysis tools can trivially prove.

\textbf{Rationale :} Prevents runaway code,
    guarantees termination,
    enables real - time deadline verification.

\textbf{STAR Compliance:
}
\textcolor{green} {
  \textbf{$\checkmark$ COMPLIANT}
}

\begin{lstlisting}[language = C]
/* COMPLIANT - Fixed iteration count */
#define MAX_RETRIES (3)

    for (uint8_t i = 0; i < MAX_RETRIES; i++) {
  if (sensor_read() == ESP_OK) {
    break;
 }
  vTaskDelay(pdMS_TO_TICKS(100));
}

/* COMPLIANT - Array with known size */
#define BQ7850_MAX_CELLS (16)

for (uint8_t cell = 0; cell < BQ7850_MAX_CELLS; cell++) {
  read_cell_voltage(cell, &voltages[cell]);
}
\end{lstlisting}

\subsubsection{Rule 3 : No Dynamic Memory After Initialization}

\textbf{Rule :
} Prohibit \texttt{malloc} /\texttt{free} after initialization phase.

\textbf{Rationale :} Dynamic allocation causes unpredictable performance,
    memory leaks, fragmentation,
    and corruption.

\textbf{STAR Compliance:
}
\textcolor{green} {
  \textbf{$\checkmark$ COMPLIANT}
}

\textbf{Best Practice :} Use pre -
    allocated static pools instead of runtime memory requests.

\begin{lstlisting}[language = C]
/* COMPLIANT - Static allocation with fixed limits */
#define MAX_ITEMS (8)
#define MAX_DESC_LEN (64)

    typedef struct {
  char items[MAX_ITEMS][MAX_DESC_LEN]; /* Static 2D array */
  uint8_t item_count;
} static_pool_t;

/* Usage */
if (pool->item_count >= MAX_ITEMS) {
  return ESP_ERR_NO_MEM;
}
strncpy(pool->items[pool->item_count], desc, MAX_DESC_LEN - 1);
pool->items[pool->item_count][MAX_DESC_LEN - 1] = '\0';
pool->item_count++;
\end{lstlisting}

\subsubsection{Rule 4 : Keep Functions Short}

\textbf{Rule :} Functions should not exceed $\sim$60 lines(
    one printed page, one statement per line)
    .

\textbf{Rationale :} Represents single verifiable logical units;
improves comprehension, review,
    and testability.

\textbf{STAR Compliance:
}
\textcolor{green} {
  \textbf{$\checkmark$ MOSTLY COMPLIANT}
}

\begin{lstlisting}[language = C]
    /* COMPLIANT - Single responsibility, short function */
    esp_err_t
    rx_motor_get_current_ma(rx_motor_handle_t *handle,
                            float *current_ma) {
  if (handle == NULL || current_ma == NULL || !handle->initialized) {
    return ESP_ERR_INVALID_ARG;
 }

  int raw_value = 0;
  esp_err_t ret = read_adc_channel(handle, &raw_value);
  if (ret != ESP_OK) {
    return ret;
 }

  *current_ma = (float)raw_value * handle->ki_propi / 1000.0F;
  return ESP_OK;
}
\end{lstlisting}

\subsubsection{Rule 5 : Use Assertions Generously}

\textbf{Rule :} Minimum two assertions per function;
assertions must be side - effect -
    free.

\textbf{Rationale :} Acts as executable documentation verifying pre -
    conditions,
    post - conditions,
    and invariants.

\textbf{STAR Compliance:
}
\textcolor{green} {
  \textbf{$\checkmark$ COMPLIANT}
}

\textbf{Note:} STAR uses explicit parameter validation instead of assertions for production code (assertions disabled in release builds).

\begin{lstlisting}[language=C]
/* COMPLIANT - Explicit validation (better than assert for production) */
esp_err_t star_pid_compute(star_pid_handle_t* handle, float setpoint,
                           float measurement, float dt, float* output)
{
  /* Pre-condition checks (effectively runtime assertions) */
  if (handle == NULL) {
    return ESP_ERR_INVALID_ARG;
 }
  if (!handle->initialized) {
    return ESP_ERR_INVALID_STATE;
 }
  if (output == NULL) {
    return ESP_ERR_INVALID_ARG;
 }
  if (dt <= 0.0F || dt > 1.0F) { /* Sanity check time delta */
    return ESP_ERR_INVALID_ARG;
 }

  /* Function logic */
  float error = setpoint - measurement;
  handle->integral += error * dt;

  /* Post-condition: integral within bounds */
  if (handle->integral > handle->integral_max) {
    handle->integral = handle->integral_max;
 } else if (handle->integral < handle->integral_min) {
    handle->integral = handle->integral_min;
 }

  *output = handle->kp * error + handle->ki * handle->integral;
  return ESP_OK;
}
\end{lstlisting}

\subsubsection{Rule 6 : Declare Data at Smallest Scope}

\textbf{Rule :} Minimize variable scope to reduce accidental corruption risks.

\textbf{Rationale :} Limits chances of accidental reference or
    corruption from other code.

\textbf{STAR Compliance:
}
\textcolor{green} {
  \textbf{$\checkmark$ COMPLIANT}
}

\begin{lstlisting}[language = C]
    /* COMPLIANT - Variables declared at minimal scope */
    void
    process_sensors(void) {
  /* Loop variable declared in for statement */
  for (uint8_t i = 0; i < SENSOR_COUNT; i++) {
    /* Temporary variable for single sensor */
    float temp_c = 0.0F;
    if (read_sensor(i, &temp_c) == ESP_OK) {
      /* Process immediately, temp_c out of scope after block */
      update_telemetry(i, temp_c);
   }
 }
}
\end{lstlisting}

\subsubsection{Rule 7 : Check All Return Values and Parameters}

\textbf{Rule :} Validate all function returns and input parameters;
treat warnings as errors.

\textbf{Rationale :} Forces explicit failure -
    condition handling rather than ignoring error codes.

\textbf{STAR Compliance:
}
\textcolor{green} {
  \textbf{$\checkmark$ COMPLIANT}
}

\begin{lstlisting}[language = C]
    /* COMPLIANT - All returns checked, all parameters validated */
    esp_err_t
    star_bus_manager_add_bus(star_bus_manager_t *manager,
                             star_bus_config_t *config) {
  /* Parameter validation */
  if (manager == NULL || config == NULL) {
    return ESP_ERR_INVALID_ARG;
 }

  /* Check return value and propagate errors */
  esp_err_t ret = star_bus_config_init(config);
  if (ret != ESP_OK) {
    ESP_LOGE(TAG, "Bus initialization failed: %s", esp_err_to_name(ret));
    return ret;
 }

  /* Mutex operation checked */
  if (xSemaphoreTake(manager->mutex, pdMS_TO_TICKS(1000)) != pdTRUE) {
    ESP_LOGE(TAG, "Failed to acquire mutex");
    return ESP_ERR_TIMEOUT;
 }

  /* Critical section */
  add_to_list(manager, config);
  xSemaphoreGive(manager->mutex);

  return ESP_OK;
}
\end{lstlisting}

\subsubsection{Rule 8 : Limit Preprocessor Use}

\textbf{Rule :} Restrict to header inclusion and simple macros;
forbid token pasting, recursive macros,
    and incomplete syntactic units.

\textbf{Rationale :} Preprocessor complexity obscures operations,
    fooling developers and static analysis tools.

\textbf{STAR Compliance:
}
\textcolor{green} {
  \textbf{$\checkmark$ COMPLIANT}
}

\begin{lstlisting}[language = C]
/* COMPLIANT - Simple constant macros */
#define MAX_RETRY_COUNT (3)
#define MOTOR_PWM_FREQ_HZ (20000)
#define TIMEOUT_MS (1000)

    /* COMPLIANT - Simple inline functions (better than complex macros) */
    static inline float
    voltage_mv_to_v(uint16_t mv) {
  return (float)mv / 1000.0F;
}

/* AVOID - Complex token-pasting macros */
// #define CREATE_HANDLER(name) handler_##name##_init()  /* Too complex */
\end{lstlisting}

\subsubsection{Rule 9 : Restrict Pointer Use}

\textbf{Rule :
} Maximum one dereferencing level(\texttt{*ptr} OK, \texttt{**ptr} forbidden);
no function pointers.

\textbf{Rationale :
} Multiple indirection and function pointers make static data -
    flow analysis impossible.

\textbf{STAR Compliance:
}
\textcolor{orange} {
  \textbf{$\triangle$ !INTENTIONAL DEVIATION}
}

\textbf{Intentional Architectural Deviation:} STAR uses function pointers for \textbf{Dependency Inversion Principle(DIP)
} -- an industry-standard pattern for testability and modularity. This is a conscious trade-off for better architecture.

\textbf{DIP Pattern(Intentional Function Pointer Use) :
}

\begin{lstlisting}[language = C]
    /* INTENTIONAL DEVIATION - DIP interface with function pointers */
    typedef struct star_error_interface {
  esp_err_t (*record_error)(void *ctx, esp_err_t error, const char *message,
                            const char *file, int line, const char *func);
  bool (*can_retry)(void *ctx);
  esp_err_t (*reset_state)(void *ctx);
  void *ctx; /* Implementation context */
} star_error_interface_t;

/* Why we keep this:
 * 1. Enables mock implementations for unit testing
 * 2. Allows swapping error handlers without recompilation
 * 3. Follows SOLID principles (Dependency Inversion)
 * 4. Industry-standard pattern (used in Linux kernel, RTOS, etc.)
 */
\end{lstlisting}

\subsubsection{Rule 10 : Compile with Maximum Warnings}

\textbf{Rule :} Enable all compiler warnings at highest pedantry;
achieve zero - warning builds with static analyzers
                   .

\textbf{Rationale :} Modern compilers catch bugs before runtime.Zero warnings,
    no excuses.

\textbf{STAR Compliance:
}
\textcolor{green} {
  \textbf{$\checkmark$ COMPLIANT}
}

\begin{lstlisting}[language = bash]
#platformio.ini - Compiler flags
    build_flags = -Wall #Enable all warnings - Wextra #Extra warnings -
                  Werror #Treat warnings as errors - Wno - unused -
                  parameter #Allow unused params in interface implementations -
                  Wno - missing - field - initializers
\end{lstlisting}

\textbf{CI / CD Enforcement :} Build fails on any compiler warning.

\subsubsection{Summary:
  STAR Compliance Matrix
}

\begin{table}[H]
\centering
\caption{NASA Power of 10 Compliance Status}
\label{tab:nasa-compliance}
\begin{tabular} { | p{1cm} | p{7cm} | p{4cm} |}
\hline
\textbf{Rule} & \textbf{Description} & \textbf{STAR Status}
\\
\hline 1 & No \texttt{goto} / recursion & \textcolor{green} {
  $\checkmark$ COMPLIANT
}
\\
\hline 2 & Fixed loop bounds & \textcolor{green} { $\checkmark$ COMPLIANT}
\\
\hline 3 & No dynamic memory & \textcolor{green} { $\checkmark$ COMPLIANT}
\\
\hline 4 & Short functions($ < $60 lines) & \textcolor{green} {
  $\checkmark$ COMPLIANT
}
\\
\hline 5 & Use assertions & \textcolor{green} { $\checkmark$ COMPLIANT}
\\
\hline 6 & Minimal variable scope & \textcolor{green} {
  $\checkmark$ COMPLIANT
}
\\
\hline 7 & Check all returns & \textcolor{green} { $\checkmark$ COMPLIANT}
\\
\hline 8 & Limit preprocessor & \textcolor{green} { $\checkmark$ COMPLIANT}
\\
\hline 9 & Restrict pointers & \textcolor{orange} {
  $\triangle$ !INTENTIONAL(DIP)
}
\\
\hline 10 & Maximum warnings & \textcolor{green} { $\checkmark$ COMPLIANT}
\\
\hline
\end{tabular}
\end{table}

\subsubsection{Defensive Coding Measures}

Beyond the Power of 10 rules, STAR implements additional defensive coding measures for electrical noise resilience in the RX72N motor controller firmware.

\textbf{Independent Watchdog Timer(IWDT) :
}
\begin{itemize}
\item 120 kHz dedicated oscillator (independent of main clock)
\item 1 second timeout, fed every 4ms from 250Hz control loop
\item Detects infinite loops, deadlocks, and system hangs
\item Logs reset cause on startup for diagnostics
\end{itemize}

\textbf{Register Guard(ESD / EMI Protection) :}
\begin{itemize}
\item Captures ``golden'' register values after initialization
\item Periodically refreshes critical registers (PORT PDR, MSTPCR)
\item Recovers from transient bit-flips caused by electrical noise
\item Tracks correction count for field diagnostics
\end{itemize}

\textbf{IRQ Digital Filtering:}
\begin{itemize}
\item Hardware noise rejection on interrupt pins
\item Configurable filter clock divisor(PCLK / 1 to PCLK / 64)
\item 3 - sample debounce prevents spurious interrupts
\end{itemize}

\begin{lstlisting}[language = C]
    /* Defensive coding in motor control loop */
    while (true) {
  /* ... control logic ... */

  rx_iwdt_feed(); /* Feed watchdog every 4ms */

  if (++guard_counter >= 250) { /* Every 1 second */
    rx_register_guard_refresh();
    guard_counter = 0;
 }
}
\end{lstlisting}

\subsubsection{References}

\begin{itemize}
\item Holzmann, Gerard J. \textit {``The Power of 10: Rules for Developing Safety-Critical Code.''
} IEEE Computer, 39(6), pp. 95-97, June 2006.
\item JPL Institutional Coding Standard for the C Programming Language (JPL DOCID D-60411)
\item MISRA C:2012 Guidelines for the Use of the C Language in Critical Systems
\end{itemize}
